% Complete Chapter 7 introduction, original PDF 91-92 / printed 81-82.
% block: finite-element-forms-chapter-heading page:91 type:chapter
\clearpage\setcounter{section}{6}
\section{\BT{Finite element differential forms}{有限元微分形式}}\label{chap:finite-element-differential-forms}

% block: finite-element-forms-hodge-problem page:91 type:ordinary-paragraph
\BT{Consider now a Lipschitz polyhedral domain $\Omega\subset\mathbb R^n$ and its $L^2$ de Rham complex \eqref{eq:l2-manifold-de-rham-complex} with domain complex \eqref{eq:manifold-de-rham-domain-complex}. For any $k$, the segment}{现在考虑 Lipschitz 多面体区域 $\Omega\subset\mathbb R^n$, 以及它的 $L^2$ de Rham 复形 \eqref{eq:l2-manifold-de-rham-complex}, 定义域复形为 \eqref{eq:manifold-de-rham-domain-complex}. 对任意 $k$, 片段}
\begin{equation}\label{eq:finite-element-forms-short-complex}
H\Lambda^{k-1}(\Omega)\xrightarrow{d}H\Lambda^k(\Omega)\xrightarrow{d}H\Lambda^{k+1}(\Omega)
\end{equation}
\BT{determines a Hodge Laplace boundary value problem involving the two exterior derivatives and their adjoints, as discussed in Section~\ref{sec:abstract-hodge-laplacian}. In Theorem~\ref{thm:exterior-codifferential-adjoints} we have identified the adjoint of $(d,H\Lambda^k)$ as $(\delta,\mathring H^*\Lambda^k)$, so the strong form of the $k$-form Hodge Laplace boundary value problem is}{确定了一个涉及两个外微分及其伴随的 Hodge Laplace 边值问题, 如 Section~\ref{sec:abstract-hodge-laplacian} 所述. 在 Theorem~\ref{thm:exterior-codifferential-adjoints} 中, 已将 $(d,H\Lambda^k)$ 的伴随确定为 $(\delta,\mathring H^*\Lambda^k)$, 因而 $k$-形式 Hodge Laplace 边值问题的强表述为}
\begin{equation}\label{eq:finite-element-forms-hodge-strong}
\begin{gathered}
(\delta d+d\delta)u=f-P_{\mathsf H}^kf\text{ in }\Omega,\\
\operatorname{tr}\star u=0,\ \operatorname{tr}\star du=0\text{ on }\partial\Omega,\qquad u\perp\mathsf H^k,
\end{gathered}
\end{equation}
\BT{where $\mathsf H^k$ are the harmonic $k$-forms on $\Omega$, determined by the differential equations $dp=0$, $\delta p=0$ and the boundary condition $\operatorname{tr}\star p=0$. In the case of three dimensions, we have interpreted this problem, and related ones, in Section~\ref{sec:hodge-r3}. In Chapter~\ref{chap:approximation-hilbert-complexes} we derived the requirements needed for a successful discretization of this problem using the mixed formulation and Galerkin's method. For this we need to construct Galerkin subspaces $V_h^{k-1}$ of $H\Lambda^{k-1}(\Omega)$ and $V_h^k$ of $H\Lambda^k(\Omega)$. The key requirements of these subspaces, besides from being efficiently implementable, are that they afford good approximation, form a subcomplex, and admit a bounded cochain projection. In that case, as shown in Chapter~\ref{chap:approximation-hilbert-complexes}, the mixed Galerkin method is consistent and stable and we obtain convergence with an optimal rate.}{其中 $\mathsf H^k$ 为 $\Omega$ 上的调和 $k$-形式, 由微分方程 $dp=0$, $\delta p=0$ 和边界条件 $\operatorname{tr}\star p=0$ 确定. 在三维情形中, 已在 Section~\ref{sec:hodge-r3} 解释了此问题及相关问题. 在 Chapter~\ref{chap:approximation-hilbert-complexes} 中, 导出了使用混合表述和 Galerkin 方法成功离散化此问题所需的条件. 为此, 需要构造 $H\Lambda^{k-1}(\Omega)$ 的 Galerkin 子空间 $V_h^{k-1}$, 以及 $H\Lambda^k(\Omega)$ 的子空间 $V_h^k$. 除能够高效实现外, 这些子空间的关键要求是提供良好的逼近, 构成子复形, 并具有有界上链投影. 在这种情况下, 如 Chapter~\ref{chap:approximation-hilbert-complexes} 所示, 混合 Galerkin 方法相容且稳定, 并且得到最优阶收敛.}

% block: finite-element-forms-construction-goal page:91 type:ordinary-paragraph
\BT{In this chapter, we shall obtain efficiently implementable Galerkin subspaces of $H\Lambda^k$ by constructing them as spaces of finite elements. The desired degree of approximation will then be obtained by using a sufficiently fine mesh and/or sufficiently high degree piecewise polynomial shape functions. The challenge is to ensure that the spaces form a subcomplex admitting a bounded cochain projection.}{本章将通过构造有限元空间, 得到可高效实现的 $H\Lambda^k$ 的 Galerkin 子空间. 随后使用足够细的网格和/或足够高次的分片多项式形函数, 就能得到所需的逼近程度. 挑战在于保证这些空间构成具有有界上链投影的子复形.}

% block: complete-trimmed-families pages:91,92 type:ordinary-paragraph
\BT{Specifically, given any simplicial triangulation $\mathcal T_h$ of $\Omega$ (in any number of dimensions $n\geq1$), any form degree $0\leq k\leq n$, and any polynomial degree $r\geq1$, we shall construct two finite element subspaces of $H\Lambda^k(\Omega)$, denoted}{具体地, 给定 $\Omega$ 的任意单纯剖分 $\mathcal T_h$ (维数 $n\geq1$ 任意), 任意形式次数 $0\leq k\leq n$, 以及任意多项式次数 $r\geq1$, 将构造 $H\Lambda^k(\Omega)$ 的两个有限元子空间, 记为}
\[
\mathcal P_r\Lambda^k(\mathcal T_h)\quad\BT{\text{and}}{\text{和}}\quad\mathcal P_r^-\Lambda^k(\mathcal T_h).
\]
\BT{The former space uses as shape functions the \emph{complete} space of polynomial $k$-forms with degree at most $r$, while the latter uses a somewhat smaller space of polynomials which, following \BookCite{29}, we refer to as \emph{trimmed}. Then we show several ways to select two such spaces, one for $(k-1)$-forms, the other for $k$-forms, to solve the Hodge Laplacian. These complete and trimmed families of finite element spaces of differential forms can be viewed as the most canonical for discretization of the de Rham complex. This is one of the important conclusions of FEEC.}{前一个空间以次数不超过 $r$ 的多项式 $k$-形式的\emph{完全}空间作为形函数空间, 后一个则使用稍小的多项式空间, 按照 \BookCite{29}, 称其为\emph{修剪}空间. 随后给出几种选取两个此类空间的方法, 一个用于 $(k-1)$-形式, 另一个用于 $k$-形式, 以求解 Hodge Laplacian. 这些完全和修剪微分形式有限元空间族, 可看作离散化 de Rham 复形最典范的空间族. 这是 FEEC 的重要结论之一.}

% block: finite-element-forms-chapter-outline page:92 type:ordinary-paragraph
\BT{The chapter is organized as follows. Before defining the complete and trimmed spaces $\mathcal P_r\Lambda^k(\mathcal T_h)$ and $\mathcal P_r^-\Lambda^k(\mathcal T_h)$ in Section~\ref{sec:finite-element-differential-forms}, we establish some properties of differential forms with polynomial coefficients in Sections~\ref{sec:complete-polynomial-forms}--\ref{sec:trimmed-polynomial-forms}. Then, in Section~\ref{sec:finite-element-space-properties}, we show how these spaces may be combined into complexes which fulfill the hypotheses required for the error analysis in Chapter~\ref{chap:approximation-hilbert-complexes}, thus establishing a full set of error bounds for our methods. In Section~\ref{sec:cubical-elements} we briefly describe analogous results for finite elements based on meshes of boxes rather than simplices. Finally, in Section~\ref{sec:historical-antecedents}, we relate the elements described in this chapter to various finite elements in the literature.}{本章安排如下. 在 Section~\ref{sec:finite-element-differential-forms} 定义完全空间和修剪空间 $\mathcal P_r\Lambda^k(\mathcal T_h)$ 与 $\mathcal P_r^-\Lambda^k(\mathcal T_h)$ 之前, 先在 Section~\ref{sec:complete-polynomial-forms}--\ref{sec:trimmed-polynomial-forms} 建立多项式系数微分形式的一些性质. 随后在 Section~\ref{sec:finite-element-space-properties}, 说明如何将这些空间组合为满足 Chapter~\ref{chap:approximation-hilbert-complexes} 误差分析所需假设的复形, 从而为我们的方法建立完整的误差界. 在 Section~\ref{sec:cubical-elements}, 简要说明基于长方体网格而非单纯形网格的有限元的类似结果. 最后在 Section~\ref{sec:historical-antecedents}, 将本章描述的单元与文献中的各种有限元联系起来.}

% block: finite-element-forms-notation page:92 type:ordinary-paragraph
\BT{Before proceeding, we make a brief remark on notation. The finite element spaces $\mathcal P_r\Lambda^k(\mathcal T_h)$ and $\mathcal P_r^-\Lambda^k(\mathcal T_h)$ which we will define in the next section consist of piecewise polynomial differential $k$-forms defined over the domain $\Omega$. We will give a unified construction of the spaces, but each space depends on the polynomial degree $r\geq0$, the form degree $k$ between $0$ and $n=\dim\Omega$, and the simplicial mesh $\mathcal T_h$. We shall also use the notation $\mathcal P_r\Lambda^k$ and $\mathcal P_r^-\Lambda^k$ for the spaces of complete and trimmed polynomial (not piecewise polynomial) differential $k$-forms of degree $r$ that are used on each simplex to construct the finite element spaces. These spaces of polynomials are defined for each $n$ and for each $r\geq0$ and $0\leq k\leq n$. Finally, we write $\mathcal P_r\Lambda^k(T)$ and $\mathcal P_r^-\Lambda^k(T)$ for the spaces of restrictions of these polynomials to a single simplex $T$. Thus our notation is analogous to a common notation for standard Lagrange finite elements of degree $r$: $\mathcal P_r(\mathcal T_h)$ for the finite element space, $\mathcal P_r$ for the space of polynomials of degree $r$, and $\mathcal P_r(T)$ for their restriction to an element.}{继续之前, 简要说明记号. 下一节将定义的有限元空间 $\mathcal P_r\Lambda^k(\mathcal T_h)$ 和 $\mathcal P_r^-\Lambda^k(\mathcal T_h)$, 由定义在区域 $\Omega$ 上的分片多项式微分 $k$-形式构成. 将对这些空间给出统一构造, 但每个空间都依赖于多项式次数 $r\geq0$, 介于 $0$ 与 $n=\dim\Omega$ 之间的形式次数 $k$, 以及单纯网格 $\mathcal T_h$. 还将用记号 $\mathcal P_r\Lambda^k$ 和 $\mathcal P_r^-\Lambda^k$ 表示次数为 $r$ 的完全和修剪多项式 (而非分片多项式) 微分 $k$-形式空间, 它们用于在每个单纯形上构造有限元空间. 这些多项式空间对每个 $n$, 每个 $r\geq0$ 及 $0\leq k\leq n$ 都有定义. 最后, 用 $\mathcal P_r\Lambda^k(T)$ 和 $\mathcal P_r^-\Lambda^k(T)$ 表示这些多项式限制在单个单纯形 $T$ 上的空间. 因此, 此记号类似于次数为 $r$ 的标准 Lagrange 有限元常用记号: $\mathcal P_r(\mathcal T_h)$ 表示有限元空间, $\mathcal P_r$ 表示次数为 $r$ 的多项式空间, $\mathcal P_r(T)$ 表示它们在一个单元上的限制.}
